% Simplified format for TDTU thesis
\usepackage{polyglossia}
\setmainlanguage{english}
\setotherlanguage{vietnamese}

\usepackage[left=3.5cm, right=2cm, top=3.5cm, bottom=3cm]{geometry}

\usepackage{graphicx,float}
\usepackage{longtable}
\usepackage{rotating}
\usepackage{pdflscape}
\usepackage{amssymb,amsmath}
\usepackage{multicol}
\usepackage{xcolor}
\usepackage{fontspec}
\usepackage{listings}
% Define custom languages BEFORE \lstset
\lstdefinelanguage{docker}{
  keywords={FROM, RUN, CMD, LABEL, EXPOSE, ENV, ADD, COPY, ENTRYPOINT, VOLUME, USER, WORKDIR, ARG, ONBUILD, HEALTHCHECK, MAINTAINER, LABEL},
  sensitive=true,
  comment=[l]{\#},
  morestring=[b]',
  morestring=[b]"
}
\lstdefinelanguage{yaml}{
  keywords={version, services, build, image, ports, environment, depends_on, networks, volumes, healthcheck, deploy, replicas, command, container_name, restart, labels, expose},
  sensitive=false,
  comment=[l]{\#},
  morestring=[b]',
  morestring=[b]"
}
\lstdefinelanguage{nginx}{
  keywords={events, http, server, location, upstream, listen, server_name, root, index, proxy_pass, proxy_set_header, worker_connections, error_log, access_log},
  sensitive=false,
  comment=[l]{\#},
  morestring=[b]',
  morestring=[b]"
}
\lstdefinelanguage{sql}{
  keywords={CREATE, TABLE, DATABASE, INSERT, SELECT, UPDATE, DELETE, FROM, WHERE, PRIMARY, KEY, FOREIGN, REFERENCES, INDEX, TIMESTAMP, DEFAULT, AUTO_INCREMENT, UNIQUE, ENUM, IF, NOT, EXISTS, USE, VALUES, IGNORE, ENGINE, CHARSET, COLLATE, FULLTEXT},
  sensitive=false,
  comment=[l]{--},
  morestring=[b]',
  morestring=[b]"
}
\lstdefinelanguage{javascript}{
  keywords={const, let, var, function, async, await, return, import, export, try, catch, if, else, while, for, class, new, try, catch, throw},
  sensitive=false,
  comment=[l]{//},
  morecomment=[s]{/*}{*/},
  morestring=[b]',
  morestring=[b]"
}
\lstset{
  basicstyle=\ttfamily\scriptsize,
  columns=fullflexible,
  xleftmargin=14pt,
  xrightmargin=2pt,
  numbers=left,
  stepnumber=1,
  frame=single,
  breaklines=true,
  breakatwhitespace=false,
  postbreak=\mbox{\textcolor{gray}{$\hookrightarrow$}\space},
  tabsize=2,
  showstringspaces=false,
  captionpos=b,
  language=tex,
  alsolanguage=docker,
  alsolanguage=bash,
  alsolanguage=sql,
  alsolanguage=yaml,
  alsolanguage=nginx,
  alsolanguage=javascript
}
\usepackage{scrextend}
\usepackage{setspace}
\usepackage{indentfirst}

% Font setup
\setmainfont{Times New Roman}

% TDTU body-text size.  The standard report class only provides 10/11/12pt
% options, so changefontsizes is used to make the effective normal text
% exactly 13pt (including lists, captions and other standard size commands).
\changefontsizes{13pt}

% TOC chapter formatting
% Add "CHAPTER" prefix to chapter entries in the Table of Contents.
% We achieve this with a clean sed post-process on thesis.toc:
%   \contentsline {chapter}{\numberline {N}Title}{page}{...}
% becomes:
%   \contentsline {chapter}{\numberline {CHAPTER N}Title}{page}{...}
% This is the most reliable approach because it touches only the .toc file.
% The actual heading on the page is already formatted as "CHAPTER N: Title"
% via the \def\@makechapterhead override below.

% Number subsubsections as x.y.z and include them in TOC
\setcounter{secnumdepth}{3}
\setcounter{tocdepth}{3}
\makeatletter
\@addtoreset{subsubsection}{section}
\makeatother

% Chapter heading format
\makeatletter
\def\@makechapterhead#1{%
  \vspace*{-5ex}%
  {\parindent \z@ \raggedright \normalfont
    \interlinepenalty\@M
    % 16pt chapter title; the second argument is kept at 16pt so that the
    % global 1.5 line-spacing setting does not make chapter headings oversized.
    \fontsize{16pt}{16pt}\selectfont
    \bfseries \color{black}
    \MakeUppercase{Chapter \thechapter: #1}\par\nobreak
    \vskip 20\p@
  }}
\makeatother

% Chapter* formatting (for unnumbered chapters like TOC, Abstract)
\makeatletter
\def\@makeschapterhead#1{%
  \vspace*{-4ex}%
  {\parindent \z@ \raggedright \normalfont
    \interlinepenalty\@M
    \fontsize{16pt}{16pt}\selectfont
    \bfseries \centering
    #1\par\nobreak
    \vskip 20\p@
  }}
\makeatother

% Section formatting
\makeatletter
\def\section{\@startsection{section}{1}{\z@}%
  {3.5ex \@plus1ex \@minus.2ex}%
  {2.3ex \@plus.2ex}%
  {\normalfont\normalsize\bfseries}}
\makeatother

% Subsection formatting
\makeatletter
\def\subsection{\@startsection{subsection}{2}{\z@}%
  {3.25ex\@plus1ex \@minus.2ex}%
  {1.5ex\@plus.2ex}%
  {\normalfont\normalsize\bfseries}}
\makeatother

% Subsubsection formatting - italic for 3-level headings (x.x.x)
\makeatletter
\def\subsubsection{\@startsection{subsubsection}{3}{\z@}%
  {3.25ex\@plus1ex \@minus.2ex}%
  {1.5ex\@plus.2ex}%
  {\normalfont\normalsize\itshape}}
\renewcommand{\thesubsubsection}{\thesection.\arabic{subsubsection}}
\makeatother

% Header / Footer
\usepackage{fancyhdr}
\pagestyle{fancy}
\fancyhf{}
\fancyhead[C]{\thepage}
% Keep the header aligned with the text block.  With the asymmetric margins
% (left 3.5cm, right 2cm), the text-block centre is intentionally to the
% right of the physical A4-page centre.
\renewcommand{\headrulewidth}{0pt}
\renewcommand{\footrulewidth}{0pt}
% \chapter sets the page style to plain on its opening page.  Mirror the
% header there so page numbers remain top-centred on every page.
\fancypagestyle{plain}{%
  \fancyhf{}%
  \fancyhead[C]{\thepage}%
  \renewcommand{\headrulewidth}{0pt}%
  \renewcommand{\footrulewidth}{0pt}%
}
\setlength{\headheight}{16pt}

% Keep citations and other hyperlinks in plain black text.
\usepackage[hidelinks]{hyperref}

% Line spacing and paragraph indent
\linespread{1.5}
\setlength{\parindent}{1.27cm} % one Word tab at the left margin
\setlength{\parskip}{0pt}

% Caption
\usepackage[labelfont=bf]{caption}
\numberwithin{figure}{chapter}
\numberwithin{table}{chapter}

% Tables: allow hyphenation-like breaks inside long identifiers so that no
% cell can run into its neighbour, and add a little vertical padding.
\usepackage{array}
\renewcommand{\arraystretch}{1.2}
\setlength{\emergencystretch}{3em}
\setlength{\LTleft}{0pt}
\setlength{\LTright}{\fill}
\setlength{\LTpre}{10pt}
\setlength{\LTpost}{10pt}
\setlength{\LTcapwidth}{\textwidth}
\captionsetup[table]{width=\textwidth,position=above,skip=6pt}
\captionsetup[figure]{width=\textwidth,position=below,skip=8pt}

% Word-like proportional column types whose fractional widths sum to 1.0 of \textwidth:
%   L{0.30} -> left-aligned paragraph column (30% of \textwidth)
%   C{0.15} -> centered paragraph column (15% of \textwidth)
%   R{0.15} -> right-aligned paragraph column (15% of \textwidth)
\newcolumntype{L}[1]{>{\raggedright\arraybackslash}p{\dimexpr #1\textwidth-2\tabcolsep-1.5\arrayrulewidth\relax}}
\newcolumntype{C}[1]{>{\centering\arraybackslash}p{\dimexpr #1\textwidth-2\tabcolsep-1.5\arrayrulewidth\relax}}
\newcolumntype{R}[1]{>{\raggedleft\arraybackslash}p{\dimexpr #1\textwidth-2\tabcolsep-1.5\arrayrulewidth\relax}}

% Emulate small caps cleanly in Times New Roman (which lacks a native OTF smcp table on macOS)
\DeclareRobustCommand{\textsc}[1]{{\fontsize{10.5pt}{13pt}\selectfont\MakeUppercase{#1}}}

% Break long snake_case, kebab-case, and path identifiers cleanly in both body and tables.
\DeclareRobustCommand{\_}{\textunderscore\allowbreak}


% TOC formatting
\renewcommand{\contentsname}{TABLE OF CONTENT}
\makeatletter
\renewcommand{\@pnumwidth}{1.55em}
\renewcommand{\@dotsep}{1.5}
\renewcommand{\l@chapter}[2]{%
  \ifnum \c@tocdepth >\m@ne
    \addpenalty{-\@highpenalty}%
    \vskip 3pt \@plus \p@
    \setlength\@tempdima{1.5em}%
    \begingroup
      \parindent \z@ \rightskip \@pnumwidth
      \parfillskip -\@pnumwidth
      \leavevmode \bfseries
      \advance\leftskip\@tempdima
      \hskip -\leftskip
      \renewcommand{\numberline}[1]{\MakeUppercase{Chapter~##1:} }%
      \MakeUppercase{#1}\nobreak
      \leaders\hbox{$\m@th \mkern \@dotsep mu\hbox{.}\mkern \@dotsep mu$}\hfill
      \nobreak\kern 2pt{\bfseries \normalcolor #2}\par
    \endgroup
  \fi}
\renewcommand{\l@section}[2]{%
  \ifnum \c@tocdepth >\m@ne
    \addpenalty{-\@highpenalty}%
    \vskip 1pt \@plus \p@
    \setlength\@tempdima{2.2em}%
    \begingroup
      \parindent \z@ \rightskip \@pnumwidth
      \parfillskip -\@pnumwidth
      \leavevmode
      \advance\leftskip\@tempdima
      \hskip -\leftskip
      #1\nobreak
      \leaders\hbox{$\m@th \mkern \@dotsep mu\hbox{.}\mkern \@dotsep mu$}\hfill
      \nobreak\kern 2pt{\normalcolor #2}\par
    \endgroup
  \fi}
\renewcommand{\l@subsection}[2]{%
  \ifnum \c@tocdepth >1\relax
    \addpenalty{-\@highpenalty}%
    \vskip 0.5pt \@plus \p@
    \setlength\@tempdima{3.2em}%
    \begingroup
      \parindent \z@ \rightskip \@pnumwidth
      \parfillskip -\@pnumwidth
      \leavevmode \itshape
      \advance\leftskip 2.2em\relax
      \advance\leftskip\@tempdima
      \hskip -\@tempdima
      #1\nobreak
      \leaders\hbox{$\m@th \mkern \@dotsep mu\hbox{.}\mkern \@dotsep mu$}\hfill
      \nobreak\kern 2pt{\normalcolor #2}\par
    \endgroup
  \fi}

% Subsubsection TOC formatting (level 3, italic)
\renewcommand{\l@subsubsection}[2]{%
  \ifnum \c@tocdepth >2\relax
    \addpenalty{-\@highpenalty}%
    \vskip 0.5pt \@plus \p@
    \setlength\@tempdima{4.0em}%
    \begingroup
      \parindent \z@ \rightskip \@pnumwidth
      \parfillskip -\@pnumwidth
      \leavevmode \itshape
      \advance\leftskip 5.4em\relax
      \advance\leftskip\@tempdima
      \hskip -\@tempdima
      #1\nobreak
      \leaders\hbox{$\m@th \mkern \@dotsep mu\hbox{.}\mkern \@dotsep mu$}\hfill
      \nobreak\kern 2pt{\normalcolor #2}\par
    \endgroup
  \fi}

% LOF/LOT formatting: prepend "Figure"/"Table" and use denser dots
\renewcommand{\l@figure}[2]{%
  \ifnum \c@tocdepth >\m@ne
    \addpenalty{-\@highpenalty}%
    \vskip 1.5pt \@plus \p@
    \setlength\@tempdima{2.5em}%
    \begingroup
      \parindent \z@ \rightskip \@pnumwidth
      \parfillskip -\@pnumwidth
      \leavevmode
      \advance\leftskip\@tempdima
      \hskip -\leftskip
      Figure~#1\nobreak
      \leaders\hbox{$\m@th \mkern 1.5mu\hbox{.}\mkern 1.5mu$}\hfill
      \nobreak\kern 2pt{\normalcolor #2}\par
    \endgroup
  \fi}
\renewcommand{\l@table}[2]{%
  \ifnum \c@tocdepth >\m@ne
    \addpenalty{-\@highpenalty}%
    \vskip 1pt \@plus \p@
    \setlength\@tempdima{2.5em}%
    \begingroup
      \parindent \z@ \rightskip \@pnumwidth
      \parfillskip -\@pnumwidth
      \leavevmode
      \advance\leftskip\@tempdima
      \hskip -\leftskip
      Table~#1\nobreak
      \leaders\hbox{$\m@th \mkern 1.5mu\hbox{.}\mkern 1.5mu$}\hfill
      \nobreak\kern 2pt{\normalcolor #2}\par
    \endgroup
  \fi}

% Custom abbrv environment for abbreviations
\newenvironment{abbrv}{
  \begin{list}{}{\setlength\itemsep{0.6em}\setlength\labelwidth{2.5cm}\setlength\leftmargin{0cm}\addtolength\leftmargin{\labelwidth}\addtolength\leftmargin{\labelsep}}
  \item[]
}{\end{list}}
\makeatother

% ============================================================
% Placeholder figures
% ------------------------------------------------------------
% \placeholderfigure{<id>}{<caption>}
% Renders a framed box that reserves the space of a figure and
% prints the caption below it.  Used for figures whose artwork is
% not yet produced, so the document still compiles and the
% numbering of figures, tables and cross-references stays correct.
% Replace a placeholder by the real \includegraphics block once
% the corresponding image file exists.
% ============================================================
\makeatletter
\newcommand{\placeholderfigure}[2]{%
  \begin{figure}[H]
    \centering
    \fbox{%
      \parbox[c][5.5cm][c]{0.88\textwidth}{%
        \centering
        \color{gray}\itshape
        [ Figure placeholder --- artwork to be inserted ]%
      }%
    }%
    \caption{#2}%
    \label{fig:#1}%
  \end{figure}%
}
\makeatother